\chapter{Data Storage Plan}
\label{ch:storage-planning}
This chapter focuses on the data decisions taken while designing the MLOps pipeline. Further this section also discusses about the physical and logical structure of the dataset showcasing how the current design decisions helps in making the optimal for data analytics and downstream meaching learning tasks too.  

% Chapters~\cref{ch:candidate-datasets}--\cref{ch:pipeline-planning} fix corpus plus DAG scope; landing and joins only~\cite{lichess_open_database}.

\section{Data Source}
\label{sec:storage-source}
The data is from Lichess platform which is an open source alternative version of chess.com for playing chess online. Lichess publishes its entire games from database for public use at \url{database.lichess.org}. It is a bulk file which can be downloaded over HTTPS. Each month a single compressed file is published alongside a SHA256 checksum for validation. This project downloads those files, verifies it, and lands it in the object store named MinIO. 

\subsection{Physical Storage}
\label{sec:storage-provenance}
The raw file downloaded from Lichess which is downloaded and stored in object storage is a compressed PGN file using Zstandard. Zstandard (zstd) was chosen by Lichess because of its balance for compression ration and decompression speed. Once decompressed, the plain-text PGN is roughly five times larger. A typical month goes from ~20 GB to ~100 GB. The raw file remains in the lake as a single immutable object with no transformation is done at this stage.

% \texttt{database.lichess.org} publishes monthly \texttt{.pgn.zst} shards plus checksum pairs; verified copies sit in MinIO with inventory in~\cref{app:lichess-open-database}~\cite{lichess_open_database}.

\subsection{Logical Storage}
\label{sec:logical-storage}
The PGN format is semi-structured file format for storage chess games with metadata. There is no enforeced schema but the structure repeats reliably. Every game is a block of following components:

\begin{enumerate}
    \item A set of headers
    \item An empty line in between
    \item The move list in Standard Algebraic Notation. Some games also contains evaluation from stockfish chess engine.
\end{enumerate}



Logically we can think of the dataset as two nested entities.
\begin{enumerate}
 \item \textbf{Game:} Fixed attributes (date, players, result, opening, etc.).
 \item \textbf{Move:} An ordered list of half-moves, each carrying a move number, SAN notation, UCI coordinates, evaluation, and clock remaining. 
\end{enumerate}

The cardinality is one game to zero-or-many moves.

\section{Ingested Data Storage}
\label{ingested_data_storage}
After extraction and loading, the raw \texttt{.zst} object is transformed into a set of Parquet tables arranged in a star schema. I follow an ELT pattern: the compressed file lands unchanged in MinIO first (load), then Apache Spark reads from object storage to parse and shape the data (transform), and the curated tables are written to MinIO as Parquet and loaded into DuckDB for analytics.


\subsection{Why Parquet?}
\begin{itemize}
    \item Fast aggregations (scans only the columns needed).
    \item High compression (columns hold similar data).
    \item Predicate push-down with Spark/DuckDB (partition pruning on year/month).
\end{itemize}

\subsection{Why MinIO?}
It keeps compute and storage separate, allows direct S3-compatible reads from Spark, DuckDB, and Feast, and the raw bucket gives a permanent, immutable copy for reprocessing.

\subsection{Logical structure (Star Schema)}
The processed data is modelled as a star schema centred on \texttt{fact\_games}. The design is shown in \cref{fig:dataset-one-logical-erd} as a crow-foot ERD.

\begin{figure}[htbp]
  \centering
  \includegraphics[width=0.95\linewidth,height=0.85\textheight,keepaspectratio]{erd}
  \caption{Star schema for processed Lichess data: a central fact table of games with dimensions for players, openings, and dates.}
  \label{fig:dataset-one-logical-erd}
\end{figure}



\subsection{Justification for the star schema}
The star schema was chosen because it directly supports the two main consumer groups in the pipeline:


\begin{table}[htbp]
    \centering
    \renewcommand{\arraystretch}{1.5}
    \begin{tabular}{|p{0.3\textwidth}|p{0.65\textwidth}|}
        \hline
        \textbf{Need} & \textbf{How the star schema meets it} \\ \hline
        ML feature engineering (Feast) & Feast works with entity keys; the fact table and dimensions can be mapped to Feast FeatureViews with simple joins on player, opening, or date keys. Point-in-time joins are straightforward. \\ \hline
        Analyst / Data Scientist notebooks & A single join from the fact table to a dimension is enough for most queries; no complex multi-table logic. \\ \hline
        Storage efficiency & Player attributes (username, title) are stored once, not repeated per game. Parquet columnar compression reduces the footprint further. \\ \hline
        Incremental monthly loads & Partitioning on year/month means new months are appended; dimensions are updated with a merge without rewriting the whole fact table. \\ \hline
        Data quality checks & Each table can have a dedicated Great Expectations suite (e.g., null checks on fact columns, uniqueness on dimension keys). \\ \hline
    \end{tabular}
\end{table}

\section{Benefits, drawbacks and Mitigation of the star schema}
The star schema model is straightforward to communicate to both engineers and data scientists because breaking the dataset into games, players, openings, and dates maps naturally to the chess domain. By using this approach, we avoid the massive data duplication that typically plagues a single, wide table. This structure also perfectly complements the write-once, read-many nature of historical chess data. While queries requiring multiple dimensions do need joins, the dimensions are relatively small, and engines like DuckDB or Spark manage these columnar joins quite efficiently. We accepted a minor bit of redundancy, such as storing the \texttt{last\_known\_elo} directly in the \texttt{dim\_player} table instead of constantly recalculating it from the fact table, simply because it makes looking up current ratings significantly faster. Finally, although evolving the schema—like adding a brand new dimension—demands careful coordination, using Parquet makes this manageable by supporting schema merging and allowing new dimensions to be introduced independently over time.

\section{Alternative structure}
\begin{table}[htp!]
    \centering
    \begin{tabular}{|p{3cm}|p{3.5cm}|p{4.5cm}|p{3.5cm}|}
        \hline
        \textbf{Approach} & \textbf{Upside} & \textbf{Downside \& Mitigation} & \textbf{Why not chosen} \\ \hline
        One Big Table & Zero joins, simple CSV & Huge duplication, brittle schema, breaks temporal correctness & Duplication and temporal fidelity risk too high \\ \hline
        Star Schema (chosen) & Familiar, fits Feast, good compression, incremental & Joins needed; mitigated by columnar scans and small dimensions & --- \\ \hline
        Snowflake & Even less redundancy in dimensions & Extra joins for minimal gain (shallow ECO hierarchy) & Complexity doesn’t pay off \\ \hline
        3NF relational & Clean updates, no anomalies & Many‑table joins hurt read speed at scale & Append‑heavy workload; read performance more important \\ \hline
        Medallion layers & Clear quality gates & Overlaps with raw/processed bucket split & Redundant; our raw bucket already provides the “bronze” layer \\ \hline
    \end{tabular}
\end{table}


\section{Extraction, Loading, and Transformation (ELT)}
I follow an ELT pattern because the heavy transformation happens after the data lands in the object store MinIO.

\begin{figure}[htbp]
  \centering
  \includegraphics[width=\linewidth,height=0.55\textheight,keepaspectratio]{elt.pdf}
  \caption{ELT pattern with object storage.}
  \label{fig:elt-flow}
\end{figure}



\subsection{Extract}
An Airflow DAG fetches each monthly \texttt{.pgn.zst} file from Lichess over HTTPS and checks it against the published SHA256 checksum. A mismatch fails the task and triggers an alert.

\subsection{Load}
The verified file lands in MinIO at \texttt{s3://lichess-raw/pgn/...} without decompression or parsing. It stays there as an immutable archive until Spark reads it.

\subsection{Transform}
PySpark streams the \texttt{.zst} object from MinIO, decompresses it in-flight, and parses games with a python-chess UDF. Headers and derived fields (e.g.\ move count) go to \texttt{fact\_games}; half-moves, engine evaluations, and clock times optionally go to \texttt{fact\_moves}. Distinct players, openings, and dates are merged into \texttt{dim\_player}, \texttt{dim\_opening}, and \texttt{dim\_date}. Curated tables are written to \texttt{s3://lichess-processed/} as Parquet, with fact tables partitioned by year and month.

\subsection{Validation}
Great Expectations runs a checkpoint on each new partition (schema, null rates, value bounds). Partitions that pass are released to Feast and analyst notebooks.

